\chapter{\texttt{agents/}: the PPO configuration}
\label{ch:agents}

\section{\texttt{agents/\_\_init\_\_.py}}

\textbf{Role:} an empty file. Its presence makes \code{agents} a regular
Python package so that \code{from . import agents} in
\srcpath{tricycle/__init__.py} works and setuptools includes the folder in the
editable install. 0 lines. Without it, \code{packages.find} skips the folder,
\code{from . import agents} raises, and the task never registers. Creating
the file is the entire fix; no reinstall is needed for it
(Chapter~\ref{ch:pyproject}).

\section{\texttt{agents/rsl\_rl\_ppo\_cfg.py}}

\textbf{Role:} the hyper-parameters of the learner. Instantiated at
configuration time (Section~\ref{sec:step_cfg}), converted to a plain dict
and handed to \code{rsl_rl}'s \code{OnPolicyRunner}. 40 lines. The meaning
of each number is in Chapter~\ref{ch:rl}.

\codepart{\srcroot/luna_hifi_tasks/tricycle/agents/rsl_rl_ppo_cfg.py}{1}{9}

\begin{explain}
\lis{1}{2} The philosophy: the learner is deliberately small so that a flat
  reward curve points at the environment, not at PPO.
\lis{4}{6} Isaac Lab develop ships \code{rsl_rl}~5.x, whose native config
  uses separate \code{actor} and \code{critic} model configs. This file uses
  the older combined \code{policy} config; Isaac Lab's
  \srcpath{handle_deprecated_rsl_rl_cfg} (in
  \srcpath{source/isaaclab_rl/isaaclab_rl/rsl_rl/utils.py}) converts it at
  start-up into an \code{RslRlMLPModelCfg} for the actor and one for the
  critic, printing a deprecation notice. The behaviour is identical.
\li{8} The configclass decorator.
\li{9} The three config classes: the runner (loop settings), the
  actor-critic network, and the PPO algorithm.
\end{explain}

\codepart{\srcroot/luna_hifi_tasks/tricycle/agents/rsl_rl_ppo_cfg.py}{12}{26}

\begin{explain}
\li{12} A configclass \ldots
\li{13} \ldots\ inheriting \code{RslRlOnPolicyRunnerCfg}, which fixes
  \code{class_name = "OnPolicyRunner"} and provides defaults such as
  \code{seed = 42}, \code{device = "cuda:0"}, \code{logger = "tensorboard"},
  \code{run_name = ""}, \code{resume = False}.
\li{14} Each iteration collects 50 steps per env before updating. At 50~Hz
  that is one second of driving per env, half an episode, so every rollout
  contains episode boundaries and the value function sees both early and
  late steps.
\li{15} Stop after 300 iterations. With 4 envs that is $300\times50\times4 =
  60{,}000$ environment steps. Override with \code{--max_iterations}.
\li{16} Save a checkpoint every 50 iterations
  (\code{model_50.pt}, \code{model_100.pt}, \ldots) plus a final one.
\li{17} The experiment name; checkpoints land in
  \srcpath{logs/rsl_rl/tricycle/<timestamp>/}.
\li{19} The network config.
\li{20} Initial standard deviation of the Gaussian policy: 1.0 in action
  units. Since actions are clamped to $[-1,1]$, early exploration is
  essentially uniform over the whole range, which is what a task with an
  unknown good action wants.
\li{21} Actor: two hidden layers of 64 units.
\li{22} Critic: the same shape.
\li{23} ELU activations for both.
\lis{24}{25} No running observation normalisation for either network. The
  observations are already in sensible ranges (velocities in m/s, angles in
  rad, actions in $[-1,1]$). These two fields replace the deprecated
  runner-level \code{empirical_normalization = False} that this file used to
  set; they are what the conversion helper would have derived from it.
\li{26} End of the policy config.
\end{explain}

\codepart{\srcroot/luna_hifi_tasks/tricycle/agents/rsl_rl_ppo_cfg.py}{27}{40}

\begin{explain}
\li{27} The PPO algorithm config (\code{class_name = "PPO"} by default).
\li{28} Weight of the critic's loss relative to the policy loss.
\li{29} Clip the value-function update as well as the policy ratio
  (PPO2 style), which stabilises the critic.
\li{30} The PPO clipping range $\epsilon = 0.2$.
\li{31} Entropy bonus coefficient 0.01: a mild push to keep exploring.
\li{32} Four passes over each rollout.
\li{33} Two mini-batches per pass: with 4 envs $\times$ 50 steps the
  mini-batch is 100 samples.
\li{34} Initial Adam learning rate $10^{-3}$.
\li{35} Adaptive schedule: the learning rate is adjusted from the measured
  KL divergence after each mini-batch (Chapter~\ref{ch:rl}).
\li{36} Discount $\gamma = 0.99$, horizon about 100 steps = one episode.
\li{37} GAE $\lambda = 0.95$.
\li{38} Target KL of 0.01 for the adaptive schedule.
\li{39} Gradient norm clipped to 1.0.
\li{40} End.
\end{explain}

\begin{isaacbox}{from cfg to runner}
\code{train_rsl_rl.py} calls \code{agent_cfg.to_dict()} and passes the dict
to \code{OnPolicyRunner(env, cfg_dict, log_dir, device)}. The runner builds
the actor and critic from the (converted) model configs, the PPO object from
\code{algorithm}, and a rollout storage sized \code{num_envs} $\times$
\code{num_steps_per_env}. The dict is also written to
\srcpath{params/agent.yaml} in the log directory, so every run records the exact
hyper-parameters it used.
\end{isaacbox}
